%=====================================================================
%  Curriculum Vitae -- Itbaan Safwan
%
%  Structure follows the grad-school CV advice:
%    - Times New Roman throughout, name at 20pt bold
%    - contact details directly under the name
%    - EDUCATION first, then RESEARCH EXPERIENCE (research-based programs)
%    - PUBLICATIONS, CONFERENCE PRESENTATIONS, TEACHING, PROFESSIONAL
%    - no skills section and no summary; skills are shown inside the bullets
%    - organisation bold left, location right; role/supervisor left, dates right
%    - 3 to 5 bullets per entry, quantified where possible
%
%  Entries are wrapped so they never split across a page break.
%  Compile with pdflatex (run twice, for hyperref).
%=====================================================================
\documentclass[11pt,letterpaper]{article}

\usepackage[T1]{fontenc}
\usepackage[utf8]{inputenc}
\usepackage{mathptmx}          % Times New Roman text and math
\usepackage[margin=1in]{geometry}
\usepackage{enumitem}
\usepackage{needspace}
\usepackage[hidelinks]{hyperref}

\hypersetup{
  colorlinks=true, linkcolor=black, urlcolor=black,
  pdftitle={Itbaan Safwan -- Curriculum Vitae},
  pdfauthor={Itbaan Safwan}
}

\pagestyle{plain}
\setlength{\parindent}{0pt}
\setlength{\parskip}{0pt}
\raggedright

%---------------------------------------------------------------------
%  Structure
%---------------------------------------------------------------------

% ALL CAPS section heading. \needspace keeps a heading from being
% stranded at the foot of a page.
\newcommand{\cvsection}[1]{%
  \Needspace*{6\baselineskip}%
  \vspace{1.0em}%
  {\bfseries\uppercase{#1}}\par
  \vspace{0.15em}\hrule\vspace{0.5em}%
}

% Keeps a whole entry (heading, sub-line and bullets) on one page.
\newenvironment{entrybox}
  {\par\vspace{0.5em}\noindent\begin{minipage}{\linewidth}}
  {\end{minipage}\par}

% Organisation flush left in bold, location flush right
\newcommand{\orgline}[2]{\textbf{#1}\hfill #2\par}

% Role or supervisor flush left, dates flush right
\newcommand{\roleline}[2]{\textit{#1}\hfill #2\par\vspace{0.15em}}

\setlist[itemize]{leftmargin=1.3em, topsep=2pt, itemsep=1.5pt,
                  parsep=0pt, label=\textbullet}

%=====================================================================
\begin{document}

%---------------------------------------------------------------------
%  Header
%---------------------------------------------------------------------
\begin{center}
  {\fontsize{20}{24}\selectfont\bfseries Itbaan Safwan}\\[0.5em]
  \href{mailto:kodvavi100@outlook.com}{kodvavi100@outlook.com} $\vert$
  \href{https://www.linkedin.com/in/itbaan-safwan-845213254}{LinkedIn} $\vert$
  \href{https://github.com/itbaans}{GitHub}
\end{center}

%---------------------------------------------------------------------
\cvsection{Education}

\begin{entrybox}
  \orgline{B.S. Computer Science}{Karachi, Pakistan}
  \roleline{Institute of Business Administration}{June 2026}
  \begin{itemize}
    \item CGPA: 3.47 / 4.00. Relevant coursework: Neural Networks and Deep Learning,
          Introduction to Natural Language Processing, Computer Vision, Linear Algebra,
          Numerical Analysis.
    \item Final year project: \textit{Visual Question Answering for Gastrointestinal Imaging with
          Multimodal Reasoning}. Advisor: Dr.\ Muhammad Atif Tahir.
  \end{itemize}
\end{entrybox}

%---------------------------------------------------------------------
\cvsection{Research Experience}

\begin{entrybox}
  \orgline{Institute of Business Administration}{Karachi, Pakistan}
  \roleline{Lead researcher \textbar\ Supervisor: Dr.\ Muhammad Atif Tahir}{2025 -- 2026}
  \begin{itemize}
    \item Led \textit{Towards Grounded GI Endoscopy VQA via Multi-Task Learning on Small VLMs} as
          first author, testing whether auxiliary visual supervision makes small vision-language
          models answer medical questions from the correct visual evidence.
    \item Built the auxiliary supervision from existing resources rather than new annotation,
          deriving Grad-CAM localisation from a domain-pretrained GastroNet-5M classifier for
          classes without expert masks, and generating visual descriptions with Gemma-27B that
          were restricted to observable features.
    \item Fine-tuned three backbones (Florence-2, Qwen3.5, InternVL3.5) jointly on localisation,
          description and VQA using QLoRA in PyTorch on a single RTX 3090 Ti, keeping training
          conditions matched so the three could be compared; only the VQA objective runs at
          inference, so the auxiliary tasks add no inference cost.
    \item Improved VQA accuracy on all three backbones (Florence-2 from 79.26\% to 82.00\%), with
          the largest gains on appearance and spatial questions and gains increasing with
          question complexity.
    \item Evaluated grounding from the model's internal representations using token-to-patch
          cosine similarity rather than segmentation outputs; Florence-2 pointing-game accuracy
          rose from 0.333 to 0.805 in-distribution and 0.142 to 0.645 out-of-distribution
          (Wilcoxon signed-rank, $p = 1.47 \times 10^{-36}$).
    \item Accepted as a first-author paper and poster presentation at the EMA4MICCAI Workshop,
          MICCAI 2026.
  \end{itemize}
\end{entrybox}

\begin{entrybox}
  \orgline{MediaEval 2025 Medico Shared Task}{Karachi, Pakistan (presented online)}
  \roleline{Team lead \textbar\ Supervisor: Dr.\ Muhammad Atif Tahir}{2025}
  \begin{itemize}
    \item Led the team submission \textit{Multi-Task Learning for Visually Grounded Reasoning in
          Gastrointestinal VQA}, covering both subtasks: answering clinical questions about
          endoscopy images, and explaining those answers in natural language.
    \item Fine-tuned one Florence-2 backbone on three objectives at once with LoRA, selected at
          inference by a task-specific prompt, so a single model served all three tasks.
    \item Constructed two of the three training sets, generating 3,344 explanations with Gemma-27B
          under a prompt that suppressed medical terminology, and assembling 4,337 text-mask pairs
          from ClipSeg pseudo-masks and expert masks in Kvasir-SEG and Kvasir-Instrument.
    \item Showed that adapter capacity trades language quality against localisation under
          VQA-only supervision, and that joint training removes the trade-off, raising IoU from
          0.2961 to 0.7403 at the same rank while BLEU also improved.
    \item Scored BLEU 0.4539, ROUGE-L 0.6531, CHrF++ 63.79 and BERTScore F1 0.9479 on the official
          private test set, earning the Distinctive Award and 2nd place in both subtasks.
  \end{itemize}
\end{entrybox}

\begin{entrybox}
  \orgline{MAHED Shared Task, ArabicNLP 2025}{Karachi, Pakistan}
  \roleline{Sole author, independent project}{2025}
  \begin{itemize}
    \item Ran an independent comparison of three multimodal fusion architectures for Arabic
          hateful meme detection, holding the backbones fixed (CLIP ViT-B/32 for the image,
          MARBERT for the Arabic caption) so that only the point of fusion varied.
    \item Implemented and trained all three systems in PyTorch: cross-attention fusion,
          progressive 1-D convolutional fusion, and a two-stage model using CLIP-style contrastive
          pre-training followed by gated classification.
    \item Handled a heavily skewed label distribution with a weighted focal loss rather than
          cross-entropy.
    \item Diagnosed an experimental error in the submitted runs, where all three models had been
          trained on the 70\% validation split; after retraining on the full labelled set,
          progressive convolutional fusion reached 0.779 macro F1, 6.1 points above the submitted
          result and 0.014 above the other two architectures.
    \item Published as a sole-authored working-notes paper in the ACL Anthology, at ArabicNLP 2025
          co-located with EMNLP 2025.
  \end{itemize}
\end{entrybox}

%---------------------------------------------------------------------
\cvsection{Publications}

\begin{entrybox}
\begin{enumerate}[leftmargin=1.6em, topsep=2pt, itemsep=5pt, parsep=0pt]
  \item \textbf{Safwan, I.}, Khan, R., Shaikh, M. A., \& Tahir, M. A. (2026). Towards Grounded GI
        Endoscopy VQA via Multi-Task Learning on Small VLMs. \textit{EMA4MICCAI Workshop, MICCAI
        2026}. (accepted, poster presentation)
        \href{https://arxiv.org/pdf/2607.27122}{[preprint]}
  \item \textbf{Safwan, I.}, Shaikh, M. A., Haaris, M., Khan, R., \& Tahir, M. A. (2025).
        Multi-Task Learning for Visually Grounded Reasoning in Gastrointestinal VQA.
        \textit{Proceedings of the MediaEval 2025 Multimedia Evaluation Workshop}, CEUR-WS.
        \href{https://arxiv.org/pdf/2511.04384}{[paper]}
  \item \textbf{Safwan, I.} (2025). Thinking Nodes at MAHED: A Comparative Study of Multimodal
        Architectures for Arabic Hateful Meme Detection. \textit{Proceedings of ArabicNLP 2025},
        712--719. Association for Computational Linguistics.
        \href{https://aclanthology.org/2025.arabicnlp-sharedtasks.98.pdf}{[paper]}
\end{enumerate}
\end{entrybox}

%---------------------------------------------------------------------
\cvsection{Conference Presentations}

\begin{entrybox}
\begin{itemize}
  \item EMA4MICCAI Workshop, MICCAI 2026. \textbf{Poster presentation}, Oct 2026.
        \textit{Towards Grounded GI Endoscopy VQA via Multi-Task Learning on Small VLMs}.
  \item MediaEval 2025 Multimedia Evaluation Workshop, Dublin and online.
        \textbf{Oral presentation}, Oct 2025. \textit{Multi-Task Learning for Visually Grounded
        Reasoning in Gastrointestinal VQA}.
\end{itemize}
\end{entrybox}

%---------------------------------------------------------------------
\cvsection{Teaching Experience}

\begin{entrybox}
  \orgline{Institute of Business Administration}{Karachi, Pakistan}
  \roleline{Teaching Assistant, Theory of Automata}{Jan 2026 -- May 2026}
  \begin{itemize}
    \item Ran weekly tutorials on the core concepts of the course, working through problem sets
          with students.
    \item Supported students through labs and office hours, and worked with faculty to design,
          evaluate and grade assignments.
  \end{itemize}
\end{entrybox}

\begin{entrybox}
  \orgline{Institute of Business Administration}{Karachi, Pakistan}
  \roleline{Teaching Assistant, Neural Networks and Deep Learning}{Aug 2025 -- Dec 2025}
  \begin{itemize}
    \item Supported students through labs and office hours, covering both the theory and its
          implementation in code.
    \item Worked with faculty to design, evaluate and grade assignments, and to bring current
          tooling into the course material.
  \end{itemize}
\end{entrybox}

%---------------------------------------------------------------------
\cvsection{AI Engineering Experience}

\begin{entrybox}
  \orgline{Cervana AI}{Karachi, Pakistan}
  \roleline{AI Research Intern}{June 2026 -- August 2026}
  \begin{itemize}
    \item Benchmarked 9 Arabic text-to-speech models on word error rate (measured with
          Whisper-Large-v3 ASR), UTMOS naturalness, F0 prosody variation and real-time factor;
          the model selected reached 0.046 WER at 0.23x real time.
    \item Fine-tuned an open-source TTS model for Urdu on an AWS GPU instance and reworked it for
          streaming, which brought end-to-end latency down from about 1000\,ms to about
          200\,ms.
    \item Wrote the serving API in FastAPI with voice cloning, API-key authentication and
          per-key quotas in SQLite, containerised it with Docker, and deployed it on an AWS EC2
          GPU instance.
  \end{itemize}
\end{entrybox}

\begin{entrybox}
  \orgline{Systems Limited}{Karachi, Pakistan}
  \roleline{AI Intern}{July 2025 -- September 2025}
  \begin{itemize}
    \item Built an agentic system for lead generation with LangGraph and the Gemini API,
          structured as a ReAct loop that updates lead state, keeps its own memory, and queries
          internal databases.
    \item Implemented a multi-tier retrieval pipeline over keyword, structured and unstructured
          sources, ordered by priority so that fewer model calls were needed per query.
    \item Prototyped a voice-to-voice agent in Streamlit using Whisper and KokoroTTS over
          real-time sockets, constrained to run on CPU only.
  \end{itemize}
\end{entrybox}

%---------------------------------------------------------------------
\clearpage
\cvsection{Leadership and Outreach}

\begin{entrybox}
  \orgline{Google Developer Groups}{Karachi, Pakistan}
  \roleline{Tech Lead for Generative AI}{2026}
  \begin{itemize}
    \item Designed and delivered a two-stage generative AI workshop, introducing the concepts from
          first principles and then running hands-on experimentation.
    \item Reached 50+ students through interactive sessions aimed at building confidence and
          practical exposure to the field.
  \end{itemize}
\end{entrybox}

\begin{entrybox}
  \orgline{IBA Computer Science and Data Science Society}{Karachi, Pakistan}
  \roleline{Module Head, Machine Learning and NLP}{2025, 2026}
  \begin{itemize}
    \item Led the Machine Learning and Natural Language Processing competition modules for 100+
          participants, writing original problem statements to reward creative and rigorous
          solutions.
    \item Curated high-difficulty, leakage-free datasets, including a geopolitical NLP challenge
          built from structured news data.
    \item Evaluated submissions on both leaderboard performance and written methodology reports.
  \end{itemize}
\end{entrybox}

\end{document}
